\begin{titlepage}

\vspace*{0.3cm}

Удирдагч:
\begin{flushright}
\makebox[4cm]{\dotfill} /\supname/
\end{flushright}

\begin{center}

\vspace*{0.5cm}
\textbf{{\large ТӨГСӨЛТИЙН АЖЛЫН ТӨЛӨВЛӨГӨӨ}}\\[0.5cm]

\textsc{\large Сэдэв: ''\ttitle''}\\[0.5cm]

\fittable{%
\begin{tabular}{|c|p{7.8cm}|c|p{2.2cm}|}
\hline
\textbf{№} & \makebox[7.8cm][c]{\textbf{Ажлын бүлэг, хэсгийн нэр}} & \textbf{Эзлэх хувь} & \textbf{Дуусах хугацаа} \\ \hline
1 & \textbf{Бүлэг 1: Асуудлын томьёолол (Удиртгал)}
\newline $\bullet$ Судалгааны үндэслэл, өнөөгийн нөхцөл байдал
\newline $\bullet$ SMART зорилго, зорилтууд, таамаглал
\newline $\bullet$ Хамрах хүрээ, хязгаарлалт
\newline $\bullet$ Ажлын шинэлэг тал, ач холбогдол & 8\% & . . . . . . . . . . . \\ \hline
2 & \textbf{Бүлэг 2: Холбогдох судалгааны тойм}
\newline $\bullet$ NLP, Sentiment Analysis-ын онолын үндэс
\newline $\bullet$ Toxic content болон Constructive criticism-ын ялгаа
\newline $\bullet$ Transformer, BERT, MN-BERT загваруудын харьцуулалт
\newline $\bullet$ Дор хаяж 5 эх сурвалжтай харьцуулсан хүснэгт
\newline $\bullet$ Авч үзээгүй талуудын дүн шинжилгээ & 10\% & . . . . . . . . . . . \\ \hline
3 & \textbf{Бүлэг 3: Аргазүй ба архитектур}
\newline $\bullet$ Судалгааны загвар, хоёр шатлалт ангиллын аргачлал
\newline $\bullet$ MN-BERT нарийн тохируулгын стратеги
\newline $\bullet$ Суурь загваруудын (LR, NB, SVM) тохиргоо
\newline $\bullet$ AI хэрэгслийн ашиглалтын ил тод байдал
\newline $\bullet$ Үнэлгээний метрик, туршилтын аргачлал & 15\% & . . . . . . . . . . . \\ \hline
4 & \textbf{Бүлэг 4: Өгөгдлийн боловсруулалт}
\newline $\bullet$ Найман нийтэд нээлттэй эх сурвалжаас өгөгдөл цуглуулалт
\newline $\bullet$ Кирилл шүүлтүүр, урьдчилсан боловсруулалтын дамжлага
\newline $\bullet$ 4 ангиллын annotation стратеги
\newline $\bullet$ Inter-annotator agreement хэмжилт
\newline $\bullet$ 70/15/15 харьцаатай train/val/test хуваалт & 15\% & . . . . . . . . . . . \\ \hline
5 & \textbf{Бүлэг 5: Загвар хөгжүүлэлт ба дамжлага}
\newline $\bullet$ MN-BERT архитектурын дэлгэрэнгүй тайлбар
\newline $\bullet$ Stage 1, Stage 2 сургалтын дамжлага
\newline $\bullet$ Гиперпараметрийн оновчлол
\newline $\bullet$ Ангиллын тэнцвэржүүлэлт (SMOTE суурь загварт / Focal Loss MN-BERT-д)
\newline $\bullet$ Gradio demo интеграц, туршилтын бүртгэл & 20\% & . . . . . . . . . . . \\ \hline
6 & \textbf{Бүлэг 6: Үр дүн ба хэлэлцүүлэг}
\newline $\bullet$ Stage 1, Stage 2 үр дүнгийн хүснэгт
\newline $\bullet$ End-to-end системийн гүйцэтгэлийн үнэлгээ
\newline $\bullet$ Алдааны дүн шинжилгээ, жишээ тайлбар
\newline $\bullet$ Суурь загвартай харьцуулсан дүгнэлт
\newline $\bullet$ Шалтгаан--үр дагаврын тайлбар & 15\% & . . . . . . . . . . . \\ \hline
7 & \textbf{Бүлэг 7: Ёс зүйн асуудал}
\newline $\bullet$ Өгөгдлийн ёс зүй, зөвшөөрөл, GDPR
\newline $\bullet$ Annotation-ы ёс зүй, Fleiss' Kappa
\newline $\bullet$ Bias ба тэгш байдлын шинжилгээ
\newline $\bullet$ Нууцлал, PII-ийн хамгаалалт
\newline $\bullet$ Хос хэрэглээний эрсдэлийн үнэлгээ & 7\% & . . . . . . . . . . . \\ \hline
8 & \textbf{Бүлэг 8: Дүгнэлт}
\newline $\bullet$ Бүлэг тус бүрийн оруулсан хувь нэмэр
\newline $\bullet$ Ш-1...Ш-7, НШ-1...НШ-5 шаардлагын биелэлтийн хүснэгт
\newline $\bullet$ Хязгаарлалтууд
\newline $\bullet$ Цаашдын судалгааны чиглэл
\newline $\bullet$ Нийгмийн нөлөөллийн үнэлгээ & 10\% & . . . . . . . . . . . \\ \hline
\multicolumn{2}{|r|}{\textbf{Нийт:}} & \textbf{100\%} & \\ \hline
\end{tabular}}

\vspace{1.5cm}
Төлөвлөгөөг боловсруулсан оюутан: \makebox[3cm]{\dotfill} /\shortname/

\end{center}

\newpage

%-------------------------------------------------------------------------------
%	REVIEW PAGE (ШҮҮМЖИЙН ХУУДАС)
%-------------------------------------------------------------------------------

\begin{center}
\textbf{{\large ТӨГСӨЛТИЙН АЖЛЫН ШҮҮМЖИЙН ХУУДАС}}\\[1cm]
\end{center}

\noindent \textbf{Оюутны нэр:} \shortname \\
\noindent \textbf{Сэдэв:} \ttitle \\
\noindent \textbf{Удирдагч:} \supname \\
\noindent \textbf{Шүүмжлэгч:} \readname \\

\vspace{0.5cm}
\noindent \textbf{Шүүмжлэгчийн дүгнэлт:}
\begin{framed}
\vspace{10cm} % Space for the reviewer to write
\end{framed}

\vspace{0.5cm}
\noindent \textbf{Үнэлгээний санал:}
\begin{table}[H]
\centering
\begin{tabular}{|l|c|c|}
\hline
\textbf{Үнэлгээний шалгуур} & \textbf{Хамгийн их оноо} & \textbf{Авсан оноо} \\ \hline
Сэдвийн ач холбогдол, шинэлэг тал & 10 & \\ \hline
Судалгааны арга зүй, архитектур & 20 & \\ \hline
Өгөгдлийн боловсруулалт, цэвэрлэгээ & 15 & \\ \hline
Загвар хөгжүүлэлт, туршилтын үр дүн & 25 & \\ \hline
Дүгнэлт, саналын үндэслэл & 10 & \\ \hline
Дипломын бичиглэл, стандартын биелэлт & 10 & \\ \hline
Асуултанд хариулсан байдал & 10 & \\ \hline
\textbf{Нийт} & \textbf{100} & \\ \hline
\end{tabular}
\end{table}

\vspace{1cm}
\begin{flushright}
Шүүмжлэгчийн гарын үсэг: \makebox[4cm]{\dotfill} \\
Огноо: . . . . . . . . . . . . . .
\end{flushright}

\end{titlepage}
